\section{Cycle variances, tuple entropy, and collisions}
\label{lifts:cycle-information}

The uniform fixed-list entropy bound \eqref{noise:uniform-entropy}
already implies the mean entropy conclusion studied here. We give two
other mechanisms for energy contraction: an exact identity comparing
one cycle's block sums with the preceding cycle's row variances, and
a collision formula for signed weights transported by actual cards.
The first yields a one-cycle recurrence and mean entropy decay; the
second yields a two-pass variation estimate. In both arguments the
input list is sampled uniformly, so available-set concentration is
available unconditionally. Their structural identities, rather than
a stronger list guarantee, are the purpose of this section.

Let $\pi_k$ be the uniform law on ordered injections of $k$ labels into
$V$, and use the relative entropy defined in Section~\ref{sec:prelim}.
An ordering of a fixed input set merely records its restriction law as
a tuple; it does not change relative entropy. Thus the mean below can
equally be viewed as an average over omitted input sets. Section~\ref{sec:c-random-domains}
will give an alternative vector proof of that information bound and
apply it through a different Fourier transfer.

\begin{proposition}[Cycle variance and tuple entropy]
\label{lifts:cycle-entropy}
For every fixed $0<\delta<1$ there is an integer $R=R(\delta)$ such that,
for a uniform ordered $k$-tuple $X$ of distinct inputs with
$k\le(1-\delta)n$, its image law $\mu_X$ after $Rd$ shuffles satisfies
\[
 \mathbb E_X D(\mu_X\|\pi_k)=o(1).
\]
Here $D$ is relative entropy with natural logarithms. More precisely,
for each rank of that random tuple, let $a_t$ be the squared
conditional next-card mass error averaged over both the input tuple
and the shuffle. Then $a_t$ satisfies, for $r\ge1$,
\begin{equation}
 a_{(r+1)d}\le \kappa_{\delta,d} a_{rd}
       +d n(n+1)e^{-\delta\sqrt n/8},\qquad
 \kappa_{\delta,d}=(1-\delta/4)^{\lfloor d/2\rfloor}
                         +2^{-\lceil d/2\rceil}.
 \label{lifts:cycle-recurrence}
\end{equation}
Thus one may choose $\alpha>0$ with $\kappa_{\delta,d}\le n^{-\alpha}$ eventually
and then any fixed $R$ with $\alpha(R-1)>4$.
\end{proposition}

\begin{proof}
At rank $j$, let $\mathcal F_t$ contain the entire realized input
tuple $X$ and the first $j-1$ cards' paths through time $t$. The
next-card law is conditioned on this field, so the initial tagged label
is fixed inside the conditional law. The expectation defining $a_t$
still averages over both $X$ and the switches. Let $V_t$ be the
available set, of size $h=n-j+1\ge\delta n$,
and let $I_t=\one_{V_t}$. Apply Lemma~\ref{lem:marginal-averaging}
with the preceding cards as blockers. Its prefix recursion gives the
conditional next-card law $p_t$, also when the complete blocker paths
are specified. Use the centered vector $f_t=p_t-I_t/h$ and mean energy
$a_t=\mathbb E\|f_t\|_2^2$ from Section~\ref{sec:first-route}.
The vector has zero sum, energy at most one, and nonincreasing energy.
The entropy chain rule and convexity give
\begin{equation}
 \mathbb E_XD(\mu_X\|\pi_k)
 \le \sum_{j=1}^k h\,\mathbb E\|f_{Rd}\|_2^2,
 \label{lifts:entropy-chain}
\end{equation}
since $D(q\|I/h)\le h\|q-I/h\|_2^2$ by $\log u\le u-1$.

We establish \eqref{lifts:cycle-recurrence}. Conditional on the complete
past before a layer, the expected entries of both $I$ and $f$ at each
pair are their old pair averages. At a pair $\{x,y\}$ the energy loss is
\[
 \tfrac12\big(f(x)^2I(y)+f(y)^2I(x)-2f(x)f(y)\big).
\]
Consider coordinate $i$ of the cycle starting at $rd$, and the blocks
$B$ obtained by varying the first $i-1$ coordinates. Their size is
$m=2^{i-1}$. Write $W_B=\sum_Bf_{rd}$ and $H_B=\sum_BI_{rd}$.
During the first $i-1$ stages these blocks evolve independently,
conditional on the cycle start. Just before stage $i$, their expected
entries are $W_B/m$ and $H_B/m$. If every $H_B\ge\delta m/2$,
independence of paired blocks and the loss formula give contraction
$1-\delta/4$ plus a cross term at most
$(2m)^{-1}\sum_BW_B^2$. If this count condition fails, use only
nonincrease and energy at most one. Thus
\begin{equation}
 a_{rd+i}\le(1-\delta/4)a_{rd+i-1}
 +\frac1{2m}\mathbb E\sum_BW_B^2
 +\frac\delta4\Pr\{\text{some }H_B<\delta m/2\}.
 \label{lifts:cycle-stage}
\end{equation}

The key estimate is $\mathbb E\sum_BW_B^2\le a_{rd}$ for $r\ge1$.
To see it, condition at $s=(r-1)d+i-1$ in the preceding cycle.
The remaining layers evolve independently in rows $C$ fixing the
first $i-1$ bits, each of size $q=n/m$. Let $R_C$ be the signed row
total at time $s$. At the end of that cycle each entry of row $C$
has conditional mean $R_C/q$. These are the transverse partitions
of Figure~\ref{cycles:transverse-figure}. A block varies the first $i-1$ bits,
whereas a row fixes those bits, so each block meets each row once.
Consequently the conditional mean of every $W_B$ is
$q^{-1}\sum_CR_C=0$. Its conditional variance is the sum of the
variances of its entries, because those entries lie in independent
rows. Summing these variances over the $q$ blocks gives
\begin{equation}
 \mathbb E\left[\sum_BW_B^2\mid\mathcal F_s\right]
 =\mathbb E[\|f_{rd}\|_2^2\mid\mathcal F_s]
                         -q^{-1}\sum_CR_C^2.
 \label{lifts:cycle-variance-identity}
\end{equation}
This proves the required inequality without assuming individual row
totals vanish.

Unconditionally the available set at the start of any cycle is a uniform
$h$-subset, by the independent random input tuple. Bernoulli sampling
with parameter $h/n$, conditioned on count $h$, represents this subset;
the conditioning event has probability at least $1/(n+1)$ because $h$
is a binomial mode. A Chernoff bound and union over the blocks give
\[
 \Pr\{\text{some }H_B<\delta m/2\}
       \le n(n+1)e^{-\delta m/8}.
\]
Apply \eqref{lifts:cycle-stage} for
$i=\lceil d/2\rceil+1,\ldots,d$, and nonincrease at earlier stages.
The sum of $(2m)^{-1}$ in this range is at most
$2^{-\lceil d/2\rceil}$. This proves the displayed recurrence.
Its additive error is smaller than every inverse power of $n$.
Since $a_d\le1$, the stated choice of $R$ gives
$a_{Rd}=O(n^{-4})$, uniformly in $j,k$. Equation~\eqref{lifts:entropy-chain}
is then $O(n^{-2})$, proving the proposition.
\end{proof}

\begin{proposition}[Two-pass transport estimate]
\label{lifts:transport-marginals}
For every fixed $0<p<1$, after a fixed number of complete coordinate
passes the average, over uniformly chosen input lists of
$\lfloor(1-p)n\rfloor$ distinct cards, of their image-law variation
distance from uniform tends to zero.
\end{proposition}
\begin{proof}
Write $x$ for the centered conditional field $f$ of the preceding
proof. It vanishes off the $h\ge pn$ available positions and is
centered there by $1/h$. We first estimate the block sums after one
pass, and then use them as the forcing term in the next pass.
A level-$j$ block varies coordinates $1,\ldots,j$ and fixes the suffix;
write $\mathcal B_j$ for this partition.

Condition on the full past at the first-pass start, fixing the available
set, signed field $x$, and energy $D=\sum_u x(u)^2$. Attach weight
$x(u)$ to the actual card at each starting site $u$, and let $L(y)$
be its transported weight at the end of the pass. Conditional on the
new blocker paths, the expectation of $L$ is the signed
transport-and-average field: Lemma~\ref{lem:marginal-averaging}
leaves each unvisited switch fair. This is a linear identity for fixed
starting weights, not a change in the conditioning that defines the
tag law. If $\bar x_B$ is the resulting field average on $B$, Jensen gives
\[
 \mathbb E\sum_{B\in\mathcal B_j}2^j\bar x_B^2
 \le 2^{-j}\mathbb E\sum_{B\in\mathcal B_j}
                       \left(\sum_{y\in B}L(y)\right)^2.
\]

The quadratic form on the right concerns two distinct actual cards,
not the conditionally independent walkers of
Lemma~\ref{sweeps:collision}. Use the last-differing-coordinate
argument of Proposition~\ref{cycles:negative-prop}, now with equality
of the output suffix as the target event. Let $a$ be the largest coordinate in
which their starting sites differ, and put $r=2^{-(d-j)}$. Their
probability of ending in the same level-$j$ block is exactly
\[
 \begin{cases}
 r,&a\le j,\\
 r(1-2^{-j}),&a>j.
 \end{cases}
\]
Before coordinate $a$ their switches are distinct, since that bit still
separates them. At coordinate $a$ they share a switch precisely when
their earlier output bits agree, in which case their new bits must be
opposite. After this stage their processed prefixes are distinct, so
all later switches are again distinct. If $a\le j$, the suffix bits
are therefore independent fair pairs, giving $r$. If $a>j$, equality
of the suffix through coordinate $a-1$ has probability
$2^{-(a-j-1)}$. Conditional on this equality, the first $j$ bits agree
with probability $2^{-j}$; this is exactly the case in which the shared
switch prevents equality at $a$. Otherwise equality there has probability
$1/2$, and later suffix equalities supply $2^{-(d-a)}$.
This gives the second formula.

In particular each off-diagonal coefficient differs from $r$ by at
most $2r2^{-j}$. The common coefficient contributes
$r\sum_{u\ne v}x(u)x(v)=-rD$, since $\sum_u x(u)=0$.
The diagonal contributes $D$, and the absolute error is at most
$2r2^{-j}(\sum_u|x(u)|)^2\le2r2^{-j}nD$. As $r=2^j/n$,
\begin{equation}
 \mathbb E\sum_{B\in\mathcal B_j}2^j\bar x_B^2\le3\,2^{-j}D.
 \label{lifts:transport-blocks}
\end{equation}

Now put $b=\lfloor d/2\rfloor$ and $\sigma=1-p/4$.
Unconditionally, the available set at the second-pass start is uniform.
The same subset estimate used above and a union bound over fewer than
$2n$ blocks show that every starting block of level at least $b$ has
available fraction at least $p/2$, except with probability at most
\[
 \varepsilon_n=2n(n+1)e^{-p2^b/8}.
\]
For a fixed second-pass start satisfying these count bounds, independent
blocks have expected entries equal to their starting field averages
and expected availabilities equal to their starting densities.
The pair-loss calculation therefore bounds the expected energy $E_j$
after $j$ stages by
\begin{equation}
 E_{j+1}\le\sigma E_j+\sum_{B\in\mathcal B_j}2^j\bar x_B^2,
 \qquad b\le j<d,
 \label{lifts:transport-stage}
\end{equation}
The squared terms use independence between the two input blocks; the
cross term is bounded by half the sum of their squared signed masses.
Average this inequality on the good starts, use
\eqref{lifts:transport-blocks} for its nonnegative forcing terms, and
bound energy by one on the exceptional event. Writing $a_t$ for mean
tag energy at a pass boundary gives the explicit recurrence
\[
 a_{t+2d}\le\bigl(\sigma^{d-b}+6\,2^{-b}\bigr)a_t+\varepsilon_n,
\]
because $\sum_{j=b}^{d-1}3\,2^{-j}\le6\,2^{-b}$.
No independence between energy and the count event is used.
For every fixed
$0<\beta<\min\{-\tfrac12\log_2\sigma,\tfrac12\}$, the coefficient
is at most $n^{-\beta}$ for sufficiently large $d$.
The uniform-subset assertion holds unconditionally at every pair start,
so a fixed number of pairs with total exponent greater than four gives
expected energy $o(n^{-4})$. Cauchy--Schwarz makes each next-card
conditional variation error $o(n^{-1})$, and Lemma~\ref{lem:sequential-tv}
sums these errors to $o(1)$.
\end{proof}

\subsection{The entropy and transport applications}
\label{lifts:cycle-applications}

We apply two Fourier transfers to these estimates. For the entropy
argument choose $R=R(1/8)$ and block length $Rd$; Pinsker and Jensen
give mean variation $o(1)$ for complements of $n/8$ inputs. For the
transport argument choose a fixed $C_0$ large enough for
Proposition~\ref{lifts:transport-marginals} at $p=1/8$, with block
length $C_0d$. For each law separately, the averaged partition choice
and normalized trimming in Section~\ref{sweeps:truncation} give
a retained probability $\nu$ within $o(1)$ of the original block law
$\mu$. If $H_i$ permutes block $i$ of the selected partition and fixes
its complement, the retained law satisfies all eight coset caps
$\nu(gH_i)\le3/[G:H_i]$. The sign and probability-replacement
conditions are then those of Section~\ref{sec:probability-completion}.

For the entropy route, the isotypic transfer
\cite[Proposition~\ref*{l-l:isotypic}]{partialFourier}, with
$b=8,u=1,B=3$, gives
\[
 \|\widehat\nu(\rho)\|_{\HS}^2\le24C_1D^{-5/8}.
\]
Thus four convolutions have nonsign Plancherel contribution at most
\[
 (24C_1)^4\sum_{\rho\ne\mathbf1,\sgn}D^{1-4(5/8)}
 =(24C_1)^4\sum_{\rho\ne\mathbf1,\sgn}D^{-3/2}=o(1).
\]
Here $C_1$ is the reciprocal-degree constant from
Section~\ref{sweeps:transfer-inputs}. The probability completion
therefore gives full-deck convergence after $4Rd$ physical shuffles.

For the transport route, use the orbit-span transfer
\cite[Proposition~\ref*{l-l:orbit-span}]{partialFourier}, again with
eight blocks and $u=1$:
\[
 \|\widehat\nu(\rho)\|_{\op}
 \le\sqrt{24C_1}\,D^{-5/16}.
\]
After $k$ convolutions the contribution at degree $D$ is at most
$(24C_1)^kD^{2-5k/8}$. Thus every fixed $k$ with $5k/8-2>0$
gives convergence, using the reciprocal-degree estimate for every
fixed positive exponent. In particular eight factors leave the summable
power $D^{-3}$ and give full-deck convergence after $8C_0d$ shuffles.

The same probability completion applies to each fixed number of
transport factors. Both block lengths are whole sweeps, so the
convolution times stated above are physical-shuffle times, uniformly
over deterministic starting decks.
